\documentclass{situationist-preprint}

\title{Cross-Checking and the Perceptual Analogy}
\msid{REL-21-b}
\msfamily{REL-21}
\msdiscipline{Philosophy of Religion}
\mstype{Companion manuscript}
\mskeywords{religious experience; Alston; mystical perception; cross-checking; discernment; epistemic imperialism}

\begin{document}
\maketitle

\begin{abstract}
Alston argued that the demand that religious experience be checked as sense perception is checked is epistemic imperialism, and that Christian mystical practice has tests of its own. REL-21-a holds that religious experience is prima facie evidence of low weight because it lacks cross-checking procedures. We examine Alston's candidate tests, the fruits of the experience and its conformity with tradition and scripture, against the definition of a cross-check. The first tests the subject rather than the content; the second is not independent of the experience it checks. The perceptual analogy holds for prima facie justification and fails for weight, and the calibration principle is not imperialist, since it asks of every practice what Alston's own tests try to supply.
\end{abstract}

\section{The candidate tests}

\citet[ch.~5]{Alston1991} argued that perceptual practice and Christian mystical practice are distinct doxastic practices, each with its own overrider system, and that to require the second to meet the checks of the first is to impose the standards of one practice on another. Among the tests internal to mystical practice he listed the fruits of an experience in the subject's life and its consistency with scripture and the tradition, tests with a long history in the literature of discernment \citep{Edwards1746}.

REL-21-a defines a cross-check as a procedure, other than a further experience of the same kind by the same subject, whose outcome can disagree with a given experience and whose disagreement counts against it (Definition 2.3); and it holds that the weight of a class of seemings beyond the minimum given by phenomenal conservatism is proportional to the confirmation of its reliability by such procedures (Principle P2).

\begin{proposition}\label{p:fruits}
The fruits test is a cross-check on the subject's condition, not on the content of the experience.
\end{proposition}

\begin{proof}
Humility, charity and peace following an experience can disagree with it, in that their absence counts against it. But what their absence counts against is the soundness of the subject: it suggests delusion, pride or illness. Their presence is equally to be expected if the experience was a benign and sincere seeming whose content is false, since the belief that one has been met by a loving God can itself produce these fruits. The test thus contributes to the fourth component of Definition 2.3 of REL-21-a, a theory of the conditions of error, and to nothing else.
\end{proof}

\begin{proposition}\label{p:trad}
Conformity with tradition is not an independent cross-check.
\end{proposition}

\begin{proof}
By Proposition 4.1 of REL-21-a, the contents of religious experiences are shaped by the traditions of those who have them. An experience that conforms to the tradition is therefore expected to conform whether or not its content is true, and its conformity is weak evidence of veridicality, as one copy of a newspaper is weak confirmation of another \citep[§265]{Wittgenstein1953}. The test can disagree with an experience, and its disagreement counts against it; but its agreement adds little.
\end{proof}

The perceptual cross-checks of agreement among observers placed in the same conditions and of successful prediction have no counterpart, for the reason Alston gives: God is not at our disposal, and no one can be placed in conditions under which God will be experienced if present. The explanation is good, but it explains the absence of the checks; it does not supply them, as Lemma 3.3 of REL-21-a states.

\section{Imperialism}

\begin{objection}
Principle P2 of REL-21-a is the imperialism Alston rejected: it measures every practice by the checks of sense perception.
\end{objection}

\begin{reply}
The principle is not about sense perception. It says of any class of seemings that their weight beyond the minimum is earned by confirmation of reliability, and it would apply to a practice with checks wholly unlike those of perception. Alston's own case concedes the point, since he looks for internal tests and appeals to them as what gives the practice its standing. The question is how much those tests confirm, and by Propositions~\ref{p:fruits} and \ref{p:trad} the answer is little. \citet{Gale1991} pressed a related disanalogy; our point is that the disanalogy concerns weight, not the existence of justification, which phenomenal conservatism secures for both practices alike (EPI-19).
\end{reply}

\section{Conclusion}

The perceptual analogy is sound at the level at which REL-21-a uses it, the existence of prima facie justification. It fails at the level of weight, because the tests of mystical practice either check the subject rather than the content or are not independent of the experience. Religious experience is evidence, and the internal tests of discernment are worth having, but they raise its weight only slightly.

\begin{thebibliography}{99}\small
\bibitem[Alston(1991)]{Alston1991} Alston, W.~P. (1991). \emph{Perceiving God: The Epistemology of Religious Experience}. Cornell University Press.
\bibitem[Edwards(1746)]{Edwards1746} Edwards, J. (1746). \emph{A Treatise Concerning Religious Affections}.
\bibitem[Gale(1991)]{Gale1991} Gale, R.~M. (1991). \emph{On the Nature and Existence of God}. Cambridge University Press.
\bibitem[Wittgenstein(1953)]{Wittgenstein1953} Wittgenstein, L. (1953). \emph{Philosophical Investigations}. Trans. G.~E.~M. Anscombe. Blackwell.
\end{thebibliography}

\end{document}
